problem:  
Let \((X,d,\mu)\) be a complete metric measure space satisfying the \(\mathrm{RCD}^*(K,N)\) condition for some \(K\in\mathbb{R}\) and \(N\in(1,\infty)\), and let \(p_t(x,y)\) be its minimal heat kernel.  
For \(x,y\in X\), define the **relative entropy between two heat distributions** at time \(t>0\),
\[
\mathcal{E}_t(x,y):=\int_X p_t(x,z)\,\log\frac{p_t(x,z)}{p_t(y,z)}\, d\mu(z).
\]
It is known that
\[
-\!4t\log p_t(x,y)\to d^2(x,y) \quad\text{as }t\downarrow 0
\]
for all \(x,y\), and in Euclidean space one can compute that \(\mathcal{E}_t(x,y)\sim d^2(x,y)/(4t)\).

**Open Problem:**  
Assume that \((X,d,\mu)\) satisfies the \(\mathrm{RCD}^*(K,N)\) condition.  

1. **Leading-order confirmation.**  
   Prove (or formulate sufficient conditions to ensure) that for all \(x,y\in X\),
   \[
   \lim_{t\to0^+}4t\,\mathcal{E}_t(x,y)=d^2(x,y).
   \]
   Even this “entropy Varadhan formula” is not yet established in general.

2. **Next-term asymptotic structure.**  
   Assuming the limit in (1) holds, investigate whether there exists a measurable function \(A(x,y)\) capturing second-order deviations:
   \[
   \mathcal{E}_t(x,y)=\frac{d^2(x,y)}{4t}+A(x,y)+o(1)\qquad (t\to0^+).
   \]
   Can \(A(x,y)\) be expressed synthetic-analytically in terms of the intrinsic Ricci curvature measure or through \(\Gamma_2\)-calculus quantities?  

3. **Smooth-manifold compatibility.**  
   When \(X\) is a smooth manifold with \(\mathrm{Ric}_N\ge K\), does \(A(x,y)\) reduce to an explicit integral of curvature along the minimal geodesic between \(x\) and \(y\), mirroring the second term in the classical Minakshisundaram–Pleijel expansion of heat kernel asymptotics?

Even formulating or bounding \(A(x,y)\) in the general \(\mathrm{RCD}^*(K,N)\) case would constitute a major step toward understanding how curvature influences the dissipation of entropy at infinitesimal scales in nonsmooth settings.

Why is it a "good" problem:  
This formulation corrects the scaling and focuses on the two most fundamental analytic layers of heat-kernel geometry: the leading order reproducing distance (Varadhan) and the next order encoding curvature.  Establishing an **entropy-based Varadhan formula** and a **second-order curvature correction** would enrich the synthetic theory of Ricci curvature by uncovering an analytic invariant derived entirely from diffusion behavior—requiring deep synthesis of heat kernel analysis, Γ₂-calculus, and Gromov–Hausdorff stability.  
The question is simple to state but profound in implications: *what does curvature look like to an infinitesimally diffusing particle in a nonsmooth space?*  Answering it would significantly advance the analytic and geometric understanding of RCD structures and how curvature emerges from heat propagation.